% ============================================================
% Commands.tex — valeurs propres au rapport
% Ces commandes alimentent les métadonnées PDF et les rappels
% de couverture. La page de garde elle-même est un PDF inclus
% (chapters/00-PageDeGarde.pdf).
% ============================================================

\newcommand{\reportAuthor} {%
  Chouaib \textsc{Saad}%
}

\newcommand{\reportTitle} {%
  Plateforme conversationnelle temps réel pour l'automatisation
  de la relation client dans les télécommunications%
}

\newcommand{\reportSubject} {%
  Rapport de projet de fin d'études — Génie logiciel et architecture logicielle%
}

\newcommand{\dateSoutenance} {%
  À PRÉCISER%
}

\newcommand{\studyDepartment} {%
  Département Informatique%
}

\newcommand{\ENETCOM} {%
  Institut Supérieur des Sciences Appliquées et de Technologie de Sousse%
}

\newcommand{\codePFE} {%
  FI-GL24-103%
}

\newcommand{\juryPresident} {%
  À PRÉCISER%
}
\newcommand{\juryPresidentDesc} {%
  Président du jury%
}

\newcommand{\juryMemberOne} {%
  Mme Maha \textsc{Khemaja}%
}
\newcommand{\juryMemberOneDesc} {%
  Encadrante pédagogique%
}

\newcommand{\juryMemberTwo} {%
  M. Ayoub \textsc{Mejri}%
}
\newcommand{\juryMemberTwoDesc} {%
  Encadrant professionnel%
}

\newcommand{\specialcell}[1]{%
  \begin{tabularx}{\textwidth}{@{}X@{}}#1\end{tabularx}%
}

%%%%%%%%%%%%%%%%%%%%%%%%%%%%%%%%%%%%%%%%%%%%%%%%%%%%%%%
% Commandes propres au rapport
%%%%%%%%%%%%%%%%%%%%%%%%%%%%%%%%%%%%%%%%%%%%%%%%%%%%%%%

% Planche de deux captures côte à côte, avec sous-légendes.
% Deux captures occupent ainsi la place d'une seule, ce qui permet
% d'en montrer davantage sans allonger le document.
%   #1 = légende de la planche   #2 = clé de la planche
%   #3 = fichier gauche  #4 = sous-légende gauche
%   #5 = fichier droit   #6 = sous-légende droite
\newcommand{\planche}[6]{%
  \begin{figure}[H]%
    \centering%
    \setlength{\fboxsep}{3pt}%
    \setlength{\fboxrule}{0.5pt}%
    \begin{subfigure}[t]{0.48\linewidth}%
      \centering\fbox{\includegraphics[width=\linewidth]{#3}}%
      \caption{#4}%
    \end{subfigure}\hfill%
    \begin{subfigure}[t]{0.48\linewidth}%
      \centering\fbox{\includegraphics[width=\linewidth]{#5}}%
      \caption{#6}%
    \end{subfigure}%
    \caption{#1}%
    \label{fig:#2}%
  \end{figure}%
}

% Planche en attente des captures. Même encombrement que la planche
% réelle, chaque case décrivant précisément la capture attendue.
\newcommand{\planchePlaceholder}[6]{%
  \begin{figure}[H]%
    \centering%
    \setlength{\fboxsep}{4pt}%
    \setlength{\fboxrule}{0.5pt}%
    \begin{subfigure}[t]{0.48\linewidth}%
      \centering%
      \fbox{\parbox[c][32mm][c]{0.93\linewidth}{\centering\scriptsize\textit{#3}}}%
      \caption{#4}%
    \end{subfigure}\hfill%
    \begin{subfigure}[t]{0.48\linewidth}%
      \centering%
      \fbox{\parbox[c][32mm][c]{0.93\linewidth}{\centering\scriptsize\textit{#5}}}%
      \caption{#6}%
    \end{subfigure}%
    \caption{#1}%
    \label{fig:#2}%
  \end{figure}%
}

% Emplacement d'un visuel non encore intégré.
%   #1 = légende de la figure
%   #2 = clé de référence (sans le préfixe fig:)
%   #3 = description précise du visuel attendu
% Le \label est placé DANS l'environnement figure, après la
% légende, afin que \ref renvoie bien au numéro de la figure.
\newcommand{\figurePlaceholder}[3]{%
  \begin{figure}[H]%
    \center%
    \setlength{\fboxsep}{10pt}%
    \setlength{\fboxrule}{0.5pt}%
    \fbox{%
      \parbox[c][45mm][c]{0.85\linewidth}{%
        \centering\small\textit{Emplacement réservé}\\[2mm]#3%
      }%
    }%
    \caption{#1}%
    \label{fig:#2}%
  \end{figure}%
}

% Figure réelle, au format visuel du modèle (cadre fin, centrée).
%   #1 = fichier image  #2 = largeur  #3 = légende  #4 = clé
\newcommand{\figureReport}[4]{%
  \begin{figure}[H]%
    \center%
    \setlength{\fboxsep}{5pt}%
    \setlength{\fboxrule}{0.5pt}%
    \fbox{%
      \includegraphics[width=#2]{#1}%
    }%
    \caption{#3}%
    \label{fig:#4}%
  \end{figure}%
}

% Figure vectorielle produite en TikZ, au format visuel du modèle.
%   #1 = légende  #2 = clé  #3 = nom du fichier dans figures/
% La figure est réduite si, et seulement si, elle dépasse la largeur
% du bloc de texte. Une figure plus étroite conserve son échelle.
\newcommand{\figureTikz}[3]{%
  \begin{figure}[!htb]%
    \center%
    \setlength{\fboxsep}{6pt}%
    \setlength{\fboxrule}{0.5pt}%
    \fbox{%
      \adjustbox{max width=\dimexpr\linewidth-18pt\relax,
                 max totalheight=0.82\textheight}{\input{figures/#3}}%
    }%
    \caption{#1}%
    \label{fig:#2}%
  \end{figure}%
}
